\section{Conclusion}
\label{sec:conclusion}

This paper builds a quantitative general equilibrium model that jointly
integrates three features of modern economies---structural monopsony,
heterogeneous households, and geoeconomic trade fragmentation---into a
single, computationally tractable framework. Solved via sequence-space
Jacobians, the model delivers three propositions with precise magnitudes
and welfare implications.

We show that monopsony endogenously flattens the Phillips curve by 40\%
relative to the competitive benchmark, raising the sacrifice ratio from 20.6
to 34.4 in our baseline calibration. We show that fragmentation amplifies
supply shocks through network reallocation of Domar weights, with a closed-form
expression for the TFP cost. And we characterize the welfare-optimal augmented
Taylor rule, which conditions on the wage markdown gap and the fragmentation
premium and achieves an 18\% welfare improvement over the best standard Taylor
rule.

Three results stand out as particularly novel. First, the countercyclical
markdown---the mechanism by which monopsony becomes \emph{more} severe in
recessions---is a previously unmodeled channel of monetary policy
non-neutrality. It implies that the output cost of fighting inflation is
\emph{asymmetric over the cycle}: recessions are more costly when they
coincide with high monopsony, which is precisely when central banks are
most likely to be tightening. Second, the fragmentation Jacobian block
establishes a tractable method for embedding time-varying network structure
into the SSJ framework, which may be of independent interest for analyzing
other forms of structural change (climate transition, digitalization).
Third, the cross-country calibration offers a structural explanation for
policy divergence that does not rely on differences in preferences or
credibility---only in market structure.

Several extensions remain for future work. The current model takes the
fragmentation process as exogenous; a natural extension would endogenize
$\phi_t$ as the outcome of a geopolitical game between blocs, potentially
generating strategic complementarities in tariff-setting. Similarly,
endogenizing monopsony---allowing $\varepsilon_t$ to respond to long-run
trends in firm entry, unionization, and non-compete agreements---would
connect the model to the three-decade trend of increasing concentration
documented by \citet{autor2020fall} and \citet{de2020rise}. Finally,
extending the model to include a government bond market and debt maturity
structure would allow analysis of the interaction between quantitative
easing and the monopsony-fragmentation channels identified here.

The post-2021 experience has forced macroeconomists to confront the
inadequacy of models built on competitive markets and integrated supply
chains. We hope this paper provides a step toward a framework that is both
theoretically rigorous and empirically adequate for the fractured world
of the 2020s.
